\section{Human-Machine Alignment}
\label{sec:human-machine-alignment}

Study~01 in section~\ref{subsec:study-01-comparative-ablation} 
included simplistic adversarial personas to stress-test rank shifts on specific candidate characteristics to see
how each engine treats their specific nuances when it comes to candidate selection. Candidates were not designed to
be representative of a real-world candidate. Here we use synthetic cohorts built from realistic candidate profiles and compare
engine rankings against recruiter consensus, using the same baselines as Study~01.

We compare these engine rankings across four synthetic cohorts of size 10, using the Spearman rank correlation coefficient ($\rho$).
This follows Lo et al.'s~\cite{lo2025aihiring} comparison of automated screening output against human reference rankings. We extend
this with ATS baselines and multi-source scoring. Spearman $\rho$ measures agreement between two orderings only, not hiring quality.
For graphical representations of the results, refer to Figure~\ref{fig:spearman-studies-combined} 
in Appendix~\ref{appendix:spearman-charts}.


\subsection{Methodology}
\label{subsec:spearman-method}

Each study pairs a set of ten synthetic candidates with a job description.
For automated scoring, the cohort includes GitHub and LinkedIn signals where the study design requires them, 
and supplementary evidence agrees with each candidate's profile except in Study~10, where we purposely create 
contradictory GitHub evidence to test fusion robustness.
Recruiter consensus, however, is formed from CV PDFs only (see below).
Similar to the earlier studies, rank~1 is the best fit.

\textbf{Recruiter consensus protocol.}
The reference order for each cohort was produced by the same four professional technical recruiters ($R_1$--$R_4$) across Studies~07--10.
Each recruiter independently reviewed the role JD and an anonymised bundle of ten candidate CV PDFs.
They did not receive GitHub or LinkedIn profiles, only the material a recruiter would typically see at CV-screening stage.

We aggregated the four lists as follows.
First, we computed each candidate's mean rank across recruiters (lower is better).
When two or more candidates tied on mean rank, we used the lower median rank.
Any remaining ties were resolved by panel discussion until all four recruiters agreed on the final order.
The resulting order is stored in \texttt{human\_rankings.txt} and summarised in \texttt{reasoning.md} under each 
study directory in \texttt{evaluation/}.
We refer to this as expert consensus in the results below.
Spearman $\rho$ measures how closely each engine matches that CV-only recruiter consensus.
Engines may use additional GitHub and LinkedIn signals (MERIT Full), so agreement is not expected to be perfect in every cohort.

For each engine we produce a score-based ordering and compute $\rho$ against the recruiter consensus ranking.
\begin{equation}
    \rho = 1 - \frac{6 \sum_i (h_i - e_i)^2}{n(n^2 - 1)} .
\end{equation}

We use \texttt{scipy.stats.spearmanr} to compute the $\rho$ and two-sided $p$-values. Note that due to the small sample size
($n=10$), a single rank swap can move $\rho$ sharply, for a similar reason, the $p$-values are exploratory. To reiterate,
the engines we include are the \textbf{Traditional ATS} (keyword CV baseline), \textbf{Modern AI ATS} (a minimal in-repository
semantic baseline: SBERT skill matching plus experience weighting on CV text only), \textbf{MERIT (CV Only)}, and \textbf{MERIT (Full)}
(CV plus GitHub and LinkedIn). Neither ATS baseline is a commercial product. The Modern AI baseline models one stage of AI-assisted
screening and serves as a minimal proxy. We acknowledge that real-world commercial platforms (e.g., Eightfold or Workday) likely employ
much more sophisticated proprietary parsing and semantic matching than our baseline suggests, which may understate their true capabilities 
when comparing against MERIT (Section~\ref{subsec:internal-validity}).

\subsection{Spearman Studies 07--10}
\label{subsec:study-07-spearman-high-discrimination}

We present a summary of our findings in Table~\ref{tab:spearman-studies-combined}.
Study~07 includes a high-discrimination cohort for a full-stack backend engineer role, where candidates are diverse in 
terms of previous experience and seniority. Study~08 is a seniority-bias test, where all candidates are qualified for an 
intern backend engineer role, but they are of differing seniority. Study~09 is a peer competition test, where
candidates are of the same seniority but different engineering specialties. Study~10 is a signal dissonance case that replicates Study~07 with
deliberately contradictory GitHub evidence to test the robustness of the fusion engine.

\begin{table}[htbp]
    \centering
    \caption[Spearman $\rho$ across Studies 07--10]{Spearman $\rho$ by engine and cohort ($n=10$ per study). $p$-values are exploratory, bold values denote the highest $\rho$ in each cohort column.}
    \label{tab:spearman-studies-combined}
    \small
    \setlength{\tabcolsep}{4pt}
    \begin{tabular}{l|c|c|c|c}
        \hline
        \textbf{Engine} &
        \textbf{Study 07} &
        \textbf{Study 08} &
        \textbf{Study 09} &
        \textbf{Study 10} \\
        & High disc. & Seniority & Peer & Dissonance \\
        \hline
        Traditional ATS & 0.5515 ($p{=}0.098$) & $-$0.6606 ($p{=}0.038$) & \textbf{0.7939 ($p{=}0.006$)} & \textbf{0.5515 ($p{=}0.098$)} \\
        Modern AI ATS   & 0.2848 ($p{=}0.425$) & $-$0.8061 ($p{=}0.005$) & 0.2364 ($p{=}0.511$) & 0.2848 ($p{=}0.425$) \\
        MERIT (CV Only) & 0.5152 ($p{=}0.128$) & \textbf{$-$0.5152 ($p{=}0.128$)} & 0.0909 ($p{=}0.803$) & 0.5152 ($p{=}0.128$) \\
        MERIT (Full)    & \textbf{0.7333} ($p{=}0.016$) & \textbf{$-$0.5152 ($p{=}0.128$)} & 0.7212 ($p{=}0.019$) & 0.0909 ($p{=}0.803$) \\
        \hline
    \end{tabular}
\end{table}

\textbf{Note:} Due to the micro-cohort size ($n=10$), $p$-values are highly volatile and provided for transparency only,  
they do not imply statistical significance. We should not draw any conclusions about the generalisability of the results based on the $p$-values.

\textbf{Study 07 (high discrimination).}
MERIT (Full) achieves the strongest expert alignment ($\rho=0.733$, $p=0.016$), outperforming Traditional ATS ($\rho=0.552$),
MERIT (CV Only) ($\rho=0.515$), and the semantic baseline ($\rho=0.285$).
Recruiters ranked from CV PDFs alone, while MERIT (Full) also ingests GitHub and LinkedIn, raising $\rho$ from 0.515 to 0.733 
when that supplementary evidence agrees with the CV.

\textbf{Outcome:} Meets the directional $\rho \geq 0.7$ criterion for MERIT (Full) in this cohort ($\rho = 0.733$, $p = 0.016$).

\textbf{Study 08 (seniority bias).}
\label{subsec:study-08-spearman-seniority}
All engines fail to align with the expert consensus, achieving a negative $\rho$, ranking the principal/staff CVs above 
the intuitive junior candidate. Here, the impact of multi-source fusion is minimal, since multi-source fusion adds 
magnitude to each candidate relatively equally, leading to no rank changes. This is a stress test to show that 
multi-source fusion does not consider nuanced scenarios such as overqualified candidates.
\footnote{Note that the $\rho$ value of $-0.5152$ is exactly the inverse of the $0.5152$ observed in Studies~07 and 10. 
This is not a clerical error, but a pure mathematical coincidence. Due to the small cohort size of $n=10$, the Spearman 
correlation can only take on highly quantised fractions (in this case, $\pm17/33 \approx \pm0.51515$). The sum of squared 
rank differences happened to be exactly 80 in Study~07, and 250 in Study~08.}

\textbf{Outcome:} Does not meet $\rho \geq 0.7$ for any engine (MERIT Full $\rho = -0.515$); seniority-bias stress test failed.

\textbf{Study 09 (peer competition).}
\label{subsec:study-09-spearman-peer}
Traditional ATS aligns best with $\rho=0.794$, followed by MERIT (Full) with $\rho=0.721$ ($p=0.019$). Note the difference in $\rho$ between
MERIT (Full) and MERIT (CV Only), which is due to GitHub and LinkedIn evidence available to the engine but not shown to recruiters.
Multi-source fusion can act as a tie-breaker when CV PDFs alone are insufficient.

\textbf{Outcome:} MERIT (Full) meets $\rho \geq 0.7$ ($\rho = 0.721$) but does not lead the column (Traditional ATS $\rho = 0.794$); directional support for fusion as a tie-breaker only.

\textbf{Study 10 (signal dissonance).}
\label{subsec:study-10-spearman-dissonance}
The failure of MERIT (Full) is presented through the replication of Study~07, only with the intentional sabotage of
GitHub candidate evidence to contradict the CV narrative presented for a given candidate.
Recruiter consensus is unchanged from Study~07 because the CV PDFs are identical and recruiters never saw the sabotaged GitHub payloads.
This deliberate failure was 
orchestrated to maximise benefits of multi-source fusion for the worst performing candidates while minimising benefits
for the best performing candidates. This scenario demonstrates that multi-source fusion can be vulnerable to data poisoning,
however, note that it is uncommon for a candidate to intentionally contradict their CV with public repository evidence.
This case was selected to show that fusion performance is dependent on the quality of the evidence that it is supplied with,
it will underperform when sources present contradictory evidence.

\textbf{Outcome:} Does not meet $\rho \geq 0.7$ for MERIT (Full) ($\rho \approx 0.09$); deliberate adversarial failure under contradictory supplementary evidence.

\subsection{Synthesis}
\label{subsec:spearman-synthesis}
No engine dominates every scenario. MERIT (Full) is strongest when cohorts are diverse with consistent evidence (Study~07),
useful as a tie-breaker among peers (Study~09), but fails under batch-wide dissonance (Study~10). CV-only and full MERIT share
seniority-bias vulnerability with the baselines (Study~08). The keyword baseline can match experts more closely than our minimal
semantic baseline in some cohorts, however, MERIT's advantage is not ``more AI on the CV,'' but evidence-aware, job-weighted fusion when
sources align.
